% Part: applied-modal-logics
% Chapter: epistemic-logic
% Section: relational-models

\documentclass[../../../include/open-logic-section]{subfiles}

\begin{document}

\olfileid{aml}{el}{rel}

\olsection{સંબંધાત્મક નિદર્શનો}

જ્ઞાનસંબંધી તર્કપદ્ધતિઓનો મૂળભૂત અર્થવિજ્ઞાનલક્ષી ખ્યાલ સામાન્ય
મોડલ તર્કપદ્ધતિઓ જેવો જ છે. સંબંધાત્મક નિદર્શમાં હજી પણ જગતોનો
ગણ અને કયા જગતમાં કયાં !!{propositional
variable}sને ``સત્ય'' ગણવા તે નક્કી કરતી નિમણૂક હોય છે. એક જ
કર્તાનો વિચાર કરતા હોઈએ, તો હંમેશની જેમ એક જ સુલભતા સંબંધ
હોય છે. પરંતુ બહુ-કર્તા જ્ઞાનસંબંધી તર્કમાં એ સંબંધની જગ્યાએ
સુલભતા સંબંધોનો ગણ આવે છે: કર્તા-સંકેતોના ગણ~$G$માંના દરેક~$a$
માટે એક સંબંધ.

\emph{સંબંધાત્મક નિદર્શ}માં જગતોનો ગણ હોય છે, જે દરેક કર્તા
માટેના એક દ્વિસ્થાની સુલભતા સંબંધ વડે જોડાયેલા હોય છે; સાથે
કયા જગતમાં કયાં !!{propositional variable}s સત્ય છે તે નક્કી
કરતી નિમણૂક હોય છે.

\begin{defn}
  બહુ-કર્તા જ્ઞાનસંબંધી ભાષાનું \emph{નિદર્શ} એ ત્રિક
  $\mModel{M} = \tuple{W, R, V}$ છે, જ્યાં
  \begin{enumerate}
  \item $W$ એ ``જગતો''નો અખાલી ગણ છે,
  \item દરેક $a \in G$ માટે ${R}_a$ એ~$W$ પરનો દ્વિસ્થાની
  સુલભતા સંબંધ છે, અને
  \item $V$ એવું વિધેય છે જે દરેક !!{propositional
    variable}~$p$ને શક્ય જગતોનો ગણ $V(p)$ આપે છે.
  \end{enumerate}
  $R_a ww'$ હોય ત્યારે કહીએ કે $w'$ એ~$w$માંથી
  \emph{કર્તા aને સુલભ} છે. $w \in V(p)$ હોય ત્યારે કહીએ કે
  $p$ એ~$w$માં \emph{સત્ય} છે.
\end{defn}

અહીં કાર્યરીતિ સામાન્ય મોડલ તર્ક જેવી જ છે; ફક્ત વધુ સુલભતા
સંબંધો ઉમેરાયા છે. આપેલા કર્તાના સુલભતા સંબંધને સામાન્ય રીતે
તેની માહિતી-સ્થિતિ વિશે કંઈક દર્શાવતો માનીએ છીએ. દાખલા તરીકે,
$R_a ww'$નો અર્થ ઘણી વાર એવો લઈએ છીએ કે $w'$ એ~$w$માં~$a$ને
મળેલી માહિતી સાથે સુસંગત છે. બીજા શબ્દોમાં, $w$માં તે
જગત~$w$ અને જગત~$w'$ વચ્ચે ભેદ કરી શકતો નથી.

\end{document}
